\chapter{The Ledgers}
% OUTLINE Ch7 "The Ledgers". The citation ledger (Kodo's second eyes);
% the intuition ledger (Vol 3's invention carried); THE ROOM LEDGER
% (this volume's own: every fixed-point instance with its house, its
% binding representation, and its non-funge note -- the non-funge law
% made a table of record). Gate: Kodo (+ second eyes on the citation
% ledger).

Three laws governed this book, and a book with three laws audits itself
with three ledgers. Every claim to a name and a date: the citation
ledger, as in every volume of this series. Every picture beside the cited
material it re-sees: the intuition ledger, carried from the volume
before. And --- new here, because this volume's own law demanded a place
to be checked --- every room met, listed with the argument it is bound
to and the reminder that it is no other argument's: the room ledger, so
that the non-funge law is not merely asserted in the sixth chapter but
tabulated in the seventh. If any row of any ledger fails to trace, the
book has broken its own law, and it asks to be read against exactly that.

\section*{The citation ledger}

\subsection*{Preface}
\begin{itemize}
\item Ken Thompson, 1984 --- ``Reflections on Trusting Trust,'' the Turing
Award lecture; the self-reproducing compiler backdoor surviving in the
fixed point of compilation.
\end{itemize}

\subsection*{Chapter 1 --- The Quine}
\begin{itemize}
\item Stephen Kleene, 1938 --- the recursion theorem; self-reference
guaranteed by a system's power to describe its own operations.
\item Douglas Hofstadter, 1979 (\emph{G\"odel, Escher, Bach}) --- the name
``quine,'' after W.\,V.\,O. Quine's use--mention discipline.
\end{itemize}

\subsection*{Chapter 2 --- The Diagonal}
\begin{itemize}
\item Georg Cantor, 1891 --- the diagonal argument; the uncountability of
the reals.
\item Bertrand Russell, 1901 --- the class of classes that do not contain
themselves.
\item Jules Richard, 1905, and the Berry sentence --- the diagonal run
unfenced in a language ranging over its own descriptions.
\end{itemize}

\subsection*{Chapter 3 --- The Room in Arithmetic}
\begin{itemize}
\item Kurt G\"odel, 1931 --- arithmetization, the diagonal lemma, the
first and second incompleteness theorems.
\end{itemize}

\subsection*{Chapter 4 --- The Unnameable}
\begin{itemize}
\item Alfred Tarski, 1933 --- the undefinability of truth in formalized
languages; the object/metalanguage stratification.
\end{itemize}

\subsection*{Chapter 5 --- The One Theorem}
\begin{itemize}
\item F.\,William Lawvere, 1969 --- ``Diagonal arguments and cartesian
closed categories''; the fixed-point theorem of which Cantor, G\"odel,
Tarski, Russell, and the halting problem are instances.
\end{itemize}

\subsection*{Chapter 6 --- The Rooms Never Merge}
\begin{itemize}
\item The standard fact of theory-relative independence --- the
undecidable sentence of one system is that system's own; strengthening
grows a fresh one --- via the logic literature's treatment.
\item The measuring instrument of this series --- cited as apparatus,
once, as an illustration and not an authority.
\end{itemize}

\section*{The intuition ledger}

Twelve labeled intuitions were handed over; each is listed with the cited
material it re-sees, and none adds a claim its neighbor lacks.

\begin{itemize}
\item[I1.] The self-photographing camera (Ch.~1) --- re-sees the quine's
existence: self-reproduction built, not mystical.
\item[I2.] The write-me-down sentence (Ch.~1) --- re-sees the body/data
kernel, the seam of use and mention.
\item[I3.] The list of every name-list (Ch.~2) --- re-sees Cantor's
diagonal: the object built to disagree.
\item[I4.] The barber and the survivable limit (Ch.~2) --- re-sees the
fenced-versus-unfenced diagonal: explode or inform.
\item[I5.] The arithmetic bureaucracy (Ch.~3) --- re-sees the G\"odel
sentence: a rulebook handed a claim about its own stamping.
\item[I6.] The unlabeled crossroads (Ch.~3) --- re-sees independence: a
map honest enough to decline one junction.
\item[I7.] The dictionary defining ``true'' (Ch.~4) --- re-sees Tarski:
the word that admits the liar.
\item[I8.] The room with no window onto itself (Ch.~4) --- re-sees the
metalanguage tower: each floor blind to its own truth.
\item[I9.] The fourth performance (Ch.~5) --- re-sees the suspicion that
one theorem underlies the four.
\item[I10.] The tesseract grown a dimension more honest (Ch.~5) ---
re-sees Lawvere as the published theorem: the object caught and named,
the image become a reference.
\item[I11.] The two houses on one blueprint (Ch.~6) --- re-sees the
non-funge law: one plan, many rooms, no door in the wall that is not
shared.
\item[I12.] The cookie cutter (Ch.~6) --- re-sees the same law from its
strongest angle: the shared plan is the evidence against one room, never
for it.
\end{itemize}

\section*{The room ledger}

This is the volume's own instrument, and it exists to make the non-funge
law auditable: every fixed-point room the book met, the house it formed
in, and the note that it is bound there and poolable with no other.

\begin{itemize}
\item \emph{The quine's kernel} --- house: any general-purpose language.
Bound to: that language's own tokens and rules. Non-funge: a quine is not
a quine in a foreign interpreter; each language's own.
\item \emph{Cantor's missing real} --- house: a proposed enumeration of
the reals. Bound to: that specific enumeration. Non-funge: the real it
misses is that list's, no other's.
\item \emph{Russell's class} --- house: unrestricted comprehension. Bound
to: that comprehension. Non-funge: its incoherence is that system's.
\item \emph{G\"odel's sentence} --- house: a consistent system that
arithmetizes its syntax. Bound to: that system. Non-funge: strengthen
the system and a new sentence forms; ``the G\"odel sentence'' is a
function of the system, not one object.
\item \emph{Tarski's unnameable} --- house: a sufficiently rich formal
language. Bound to: that language. Non-funge: its truth lives one floor
up, in its own tower; there is no universal metalanguage.
\item \emph{The instrument's residue} --- house: the apparatus of this
series. Bound to: that construction. Non-funge: it is that device's own,
resembling the others in plan and identical to none in place; not
``the'' residue of anything wider.
\end{itemize}

The reader should notice what the last row does: it puts the series' own
apparatus in the same table as G\"odel's and Tarski's rooms, under the
same binding column and the same non-funge note, with no exemption. That
is the whole point of a law --- it holds where the author would most like
it not to.

\section*{The closing line}

Every argument strong enough to describe its own describing grows a room
it cannot dissolve; the room's truth is not the point, it always appears,
it is roughly always the same size, and it can never be carried across to
another argument. Four houses proved it, one theorem underlies them, and
the series' own house takes its place on the street under the same law
that governs the rest. The mathematics belongs to the names above; the
pictures re-see it and add nothing; the rooms are each their own, and the
one this series carries is no exception --- one room, one argument, its
own. There is no \emph{the} residue. There is only this one, and that
one, and the law that keeps them apart. The books are closed.
